\section{Conclusion}

This paper presented VibSemAlign, a novel framework that advances vibration-based fault diagnosis by semantically aligning spectrographic representations of machinery vibrations with natural language descriptions of fault conditions. By leveraging vision-language models fine-tuned through a specialized vibration-semantic contrastive loss and integrated with model-agnostic meta-learning, VibSemAlign achieves robust zero-shot and few-shot performance across diverse domains, as demonstrated on benchmarks such as CWRU and Paderborn datasets. Our approach addresses longstanding challenges in predictive maintenance, including data scarcity, domain shifts due to varying operational conditions, and the need for interpretable diagnostics.

The core contributions of this work are threefold. First, we propose the vibration-semantic contrastive loss (VSCL), which effectively bridges visual signal patterns with textual semantics, enabling the model to discern subtle fault signatures like bearing defect modulations without extensive retraining. Second, the incorporation of cross-attention mechanisms and MAML facilitates rapid adaptation to target domains with minimal labeled samples, yielding accuracy improvements of up to 17.2 percentage points over domain adaptation baselines in zero-shot cross-domain settings. Third, the framework provides interpretable reasoning grounded in semantic prompts, offering insights into fault mechanisms---such as BPFO harmonic tracking and wear progression patterns---that surpass traditional signal processing techniques.

Looking ahead, several avenues warrant further exploration. Extending VibSemAlign to multimodal fusion---incorporating acoustic or current signals---could enhance diagnostic granularity for complex machinery. Real-time deployment on edge devices remains a priority; our INT8 quantization experiments on Jetson Nano (45\,ms latency) indicate feasibility for 22\,Hz monitoring. Additionally, scaling to larger real-world industrial datasets and exploring unsupervised domain adaptation with dynamic prompt generation would broaden applicability across diverse manufacturing environments.

In broader terms, VibSemAlign establishes a new paradigm for semantic predictive maintenance, demonstrating that vision-language models can serve as powerful reasoning engines for industrial fault diagnosis when appropriately adapted through vibration-semantic alignment, with implications for reducing unplanned downtime and improving equipment reliability in manufacturing, energy, and transportation sectors.
